\documentclass{article}
\usepackage[T1]{fontenc}
\usepackage[spanish]{babel}
\usepackage{amssymb}
\usepackage{amsmath}
\usepackage{booktabs}
\usepackage[margin=2.5cm]{geometry}
\title{Cálculo del equilibrio: Newton conjunto con jacobiano exacto\\ y precios sombra en celdas sin oferta}
\author{Notas de apoyo a la tesis (versión de Claude)}
\date{Noviembre 2026}

% Explica la subsección "Cálculo del equilibrio" de tesis_claude/tesis.tex.
% Código de referencia: claude/niu/modelo_tesis/equilibrio.py y utils.py (jacobian).

\begin{document}
\maketitle

\section{El problema que se resuelve}

Para un régimen de precios de reparación $R$ fijo, el equilibrio con dos tipos (Teorema 1 y
sección de consumidores heterogéneos) consiste en encontrar
$$z=(EV_1,\,EV_2,\,P)\in\mathbb{R}^{2n+J(\bar a-1)W}$$
tal que $F(z)=0$, con
$$F(z)=\begin{pmatrix}
EV_1-\Gamma_1(EV_1;P,R)\\[2pt]
EV_2-\Gamma_2(EV_2;P,R)\\[2pt]
\log D(P,EV_1,EV_2)-\log S(P,EV_1,EV_2)
\end{pmatrix}.$$
Los dos primeros bloques son las ecuaciones de Bellman de cada tipo ($n$ ecuaciones cada uno) y
el tercero es el vaciado de mercado celda por celda ($J(\bar a-1)W$ ecuaciones). La distribución
estacionaria $q_\tau$ no es incógnita: dados $(EV_\tau,P)$ se calculan las probabilidades de
elección, con ellas $\Omega_\tau$ y $A_\tau$, y $q_\tau$ se obtiene por recursión en
la edad. Por eso $D$ y $S$ dependen de $EV$ además de $P$:
$$D=\sum_\tau f_\tau D_\tau\big(P,q_\tau(EV_\tau,P)\big),\qquad
S=\sum_\tau f_\tau S_\tau\big(P,q_\tau(EV_\tau,P)\big).$$
Con $J=2$, $\bar a=25$ y $W=71$: $n=3{,}411$, hay $3{,}408$ precios y $10{,}230$ incógnitas.

\paragraph{Por qué $\log D-\log S$ y no $D-S$.} Las masas por celda van de $10^{-2}$ a $10^{-12}$.
Con $D-S$, una celda con masa $10^{-12}$ ya cumple cualquier tolerancia aunque su demanda sea
el doble de su oferta. En logaritmos el residuo es un error relativo, igual de exigente en todas
las celdas, y la tolerancia $\max|F|<10^{-9}$ tiene el mismo significado en todos los bloques.

\section{El método de Newton}

Newton linealiza $F$ alrededor del punto actual y resuelve el sistema lineal:
$$F(z_k)+F_z(z_k)\,\Delta z=0\quad\Longrightarrow\quad
z_{k+1}=z_k-F_z(z_k)^{-1}F(z_k),$$
donde $F_z$ es la matriz jacobiana de $10{,}230\times10{,}230$. Cerca de la solución, si $F_z$ es
invertible ahí, la convergencia es cuadrática: $\|z_{k+1}-z^*\|\le C\|z_k-z^*\|^2$. En la práctica,
el número de dígitos correctos se duplica en cada paso ($10^{-2}\to10^{-4}\to10^{-8}\to10^{-16}$), así que
desde un buen punto de arranque bastan pocas iteraciones.

\paragraph{Relación con Newton--Kantorovich.} En la sección de Bellman de la tesis, Newton--Kantorovich
es exactamente Newton aplicado solo al primer bloque con $P$ fijo:
$EV_{k+1}=EV_k-[\mathbb I-\Gamma'(EV_k)]^{-1}(EV_k-\Gamma(EV_k))$. El Newton conjunto hace lo mismo
pero mueve $EV$ y $P$ a la vez, usando todos los bloques del jacobiano.

\paragraph{Por qué conjunto y no anidado.} La alternativa natural es anidar: para cada $P$ resolver
$EV(P)$ y luego buscar $P$ con $ED(P,EV(P))=0$. El lazo externo necesita
$\mathrm dEV/\mathrm dP=-[\mathbb I-\Gamma_{EV}]^{-1}\Gamma_P$, es decir, las mismas derivadas del
jacobiano conjunto. Si en cambio el lazo externo es un tanteo (subir el precio donde hay exceso
de demanda), la convergencia es lineal y lenta, porque los precios de las $3{,}408$ celdas están muy
acoplados (un coche $(j,a,w)$ compite con sus vecinos en $a$ y $w$). Newton conjunto usa todo el
acoplamiento de una vez y no requiere resolver $EV$ con precisión en cada paso.

\section{La estructura del jacobiano}

Ordenando las incógnitas como $(EV_1,EV_2,P)$,
$$F_z=\begin{pmatrix}
\mathbb I-\Gamma_{1,EV} & 0 & -\Gamma_{1,P}\\
0 & \mathbb I-\Gamma_{2,EV} & -\Gamma_{2,P}\\
G_{EV_1} & G_{EV_2} & G_P
\end{pmatrix},\qquad G=\log D-\log S.$$
Los ceros aparecen porque el Bellman de un tipo no depende del $EV$ del otro: los tipos solo se
conectan a través de los precios. Los bloques de Bellman tienen forma cerrada.

\subsection{Bloque $\Gamma_{EV}$: la matriz de transición}

El resultado clave es la propiedad de Williams--Daly--Zachary del log-sum: la derivada de
$I(\lambda,\chi_1,\dots,\chi_m)$ respecto a $\chi_k$ es la probabilidad de escoger $k$ dentro del nido,
$$\frac{\partial I}{\partial\chi_k}=\frac{\exp(\chi_k/\lambda)}{\sum_l\exp(\chi_l/\lambda)}.$$
Aplicando la regla de la cadena a lo largo del árbol, la derivada de $EV(x)$ respecto al valor de
una hoja $d$ es la probabilidad incondicional de esa hoja, $\Pi(d\mid x)$. $EV$ entra a las hojas
solo a través de $\beta EV(\emptyset)$ y de los $v_r(h)$:
$$\frac{\partial v_r(h)}{\partial EV(x')}=\beta\Big[s(h)\,\mathbb 1\{x'=(j,\bar a)\}+(1-s(h))f(w'\mid j,a,w,r)\,\mathbb 1\{x'=(j,a+1,w')\}\Big].$$
Juntando,
$$\frac{\partial\Gamma(x)}{\partial EV(x')}=\beta\sum_{h}\Pi(h\mid x)\sum_r\Pi(r\mid h)\,[QA]_{h,x'}
=\beta\,[\Omega\,Q\,A]_{x,x'}.$$
Es decir, $\Gamma_{EV}=\beta\,\Omega QA$, como dice la tesis. Como $\Omega QA$ es estocástica,
$\|\Gamma_{EV}\|_\infty=\beta<1$ y el bloque $\mathbb I-\beta\,\Omega QA$ siempre es invertible
(es la misma razón por la que $\Gamma$ es contracción).

\subsection{Bloque $\Gamma_P$: el exceso de oferta individual}

Con $EV$ fijo, $P$ entra a las hojas en dos lugares: en el valor de vender el coche propio,
$\nu_0(x)=\mu_\tau P_{x}-T_s$, y en el de comprar un usado, $B_r(h)=\dots-\mu_\tau P_h$. Por la misma
propiedad del log-sum, para un usado $h$,
$$\frac{\partial\Gamma(x)}{\partial P_h}=\mu_\tau\Big[\mathbb 1\{h=x\}\,\big(1-\Pi(K\mid x)\big)\big(1-\Pi(D=1\mid x)\big)
-\Pi(\text{comprar }h\mid x)\Big].$$
El primer término es la probabilidad de que el dueño de $x$ lo venda (justo el factor de $S_x$) y el
segundo la de que compre $h$. Así, la fila $x$ de $\Gamma_P$ es $\mu_\tau$ por el exceso de oferta
esperado del hogar en $x$: subir el precio de un coche beneficia a quien lo vende y perjudica a quien
lo compra, en proporción a la probabilidad de cada cosa. La derivada no tiene términos de
$\partial\Pi/\partial P$ por el teorema de la envolvente.

\subsection{Bloque de mercado $G$: el difícil}

$G$ depende de $(EV,P)$ por dos vías: las probabilidades de elección (directamente en $D$ y $S$) y
la distribución $q_\tau$. Las derivadas de las probabilidades sí tienen todos los términos de un
logit anidado. Por ejemplo, en el nido de $y$,
$\partial\Pi(y\mid i,b)/\partial B(y')=\Pi(y\mid i,b)\big(\mathbb 1\{y=y'\}-\Pi(y'\mid i,b)\big)/\sigma_w$,
y esto se multiplica a lo largo de los niveles del árbol. La parte de $q$ sale de diferenciar la
condición de estacionariedad:
$$q=q\,\Omega QA,\ q\mathbf 1=1\quad\Longrightarrow\quad\mathrm dq\,(\mathbb I-\Omega QA)=q\,\mathrm d(\Omega QA),\ \ \mathrm dq\,\mathbf 1=0,$$
donde $\mathrm d(\Omega QA)=\mathrm d\Omega\,QA+\Omega Q\,\mathrm dA$ a su vez depende de las derivadas de todas las probabilidades de elección.
Escribir esto a mano es posible pero largo y propenso a errores.

\section{Cómo se calcula el jacobiano en el código}

En \texttt{modelo\_tesis/equilibrio.py} el jacobiano no se programa a mano: se calcula con
\textbf{diferenciación automática en modo directo} (\texttt{utils.jacobian}). La idea:
\begin{enumerate}
  \item \texttt{jax.linearize} evalúa $F(z)$ una vez y guarda la linealización, una función
  $v\mapsto F_z(z)\,v$ (producto jacobiano-vector) que aplica la regla de la cadena operación por
  operación sobre el mismo código que calcula $F$.
  \item Evaluando ese producto en los vectores canónicos $e_1,\dots,e_{10230}$ se obtiene cada
  columna de $F_z$. Se hace por bloques de 256 columnas (\texttt{jac\_chunk}) vectorizados en GPU,
  así que son 40 bloques y la memoria es $O(256\times10{,}230)$ por bloque.
  \item El paso es $\Delta z=-F_z^{-1}F$ con una factorización densa
  (\texttt{jnp.linalg.solve}). La matriz ocupa $10{,}230^2\times8$ bytes $\approx 0.84$ GB y la factorización cuesta
  del orden de $10^{12}$ operaciones, por eso se hace en GPU.
\end{enumerate}
Este jacobiano es \textbf{exacto}: coincide con las fórmulas analíticas de la sección anterior hasta
el error de redondeo de la máquina (el código usa doble precisión). No es una aproximación por
diferencias finitas. La diferencia importa: con diferencias finitas
$[F(z+he_i)-F(z)]/h$ el error es del orden $h+\varepsilon_{\text{máq}}/h\approx10^{-8}$, y un Newton
con ese jacobiano se estanca justo alrededor de la tolerancia de $10^{-9}$ que pedimos.

\paragraph{Sobre la palabra ``analítico'' en la tesis.} La tesis dice ``jacobiano analítico''. Es
correcto en el sentido de que son las derivadas exactas de las fórmulas del modelo, pero el lector
podría entender que están programadas a mano. Una redacción más precisa sería: ``con el jacobiano
exacto, obtenido por diferenciación automática (los bloques de Bellman coinciden con las expresiones
cerradas $\Gamma_{EV}=\beta\,\Omega QA$ y la derivada de $\Gamma$ respecto a $P$ es el exceso de oferta
individual)''.

\section{Lo que hace robusto al solver}

Newton puro solo converge si arranca cerca de la solución. El código (\texttt{solve}) agrega:
\begin{description}
  \item[Arranque.] Precios que caen 13\% al año desde el precio del nuevo, con piso en el de
  chatarra e iguales en $w$: $P_{jaw}=\max(\overline P_j\,0.87^a,\underline P_j)$. Dado ese $P$,
  500 aproximaciones sucesivas para cada $EV_\tau$.
  \item[Tope al paso en precios.] Si el paso propuesto cambia algún precio en más de 150 mil DKK,
  se reescala todo el paso. Evita que una linealización mala mande precios a valores absurdos.
  \item[Búsqueda de línea.] Se prueba $z+\lambda\Delta z$ con $\lambda=1,\tfrac12,\tfrac14,\dots$
  (hasta 25 veces) y se acepta el primero que baje $\|F\|_2$ y sea finito. Lejos de la solución
  da pasos cortos y seguros; cerca, $\lambda=1$ y la convergencia vuelve a ser cuadrática.
  \item[Rescate.] Si ningún $\lambda$ baja el residuo, se dejan fijos los precios y se vuelve a
  resolver $EV$ con aproximaciones sucesivas antes de intentar otro paso de Newton.
  \item[Criterio de paro.] $\max|F|<10^{-9}$: en Bellman, un error de $10^{-9}$ utils; en
  mercado, un error relativo de $10^{-9}$ entre demanda y oferta en cada celda.
  \item[Arranque en caliente entre regímenes.] Primero se resuelve el régimen central
  ($\zeta=0$) desde cero, y los otros arrancan desde la solución de su vecino.
\end{description}

\paragraph{En la estimación.} Cada evaluación de la verosimilitud en un $\theta$ nuevo necesita el
equilibrio. Ahí (\texttt{solve\_chord}) se reutiliza la inversa del jacobiano de un $\theta$ cercano
(Newton ``de cuerda'': cada paso cuesta una evaluación de $F$ y un producto matriz-vector), partiendo
del predictor $z+\tfrac{\mathrm dz}{\mathrm d\theta}\Delta\theta$, con
$\tfrac{\mathrm dz}{\mathrm d\theta}=-F_z^{-1}F_\theta$ por el teorema de la función implícita. Si la
cuerda deja de reducir el residuo, se refactoriza. La misma $F_z^{-1}$ da las derivadas de la
verosimilitud respecto a $\theta$ sin diferenciar a través de las iteraciones del solver.

\section{Precios sombra en celdas casi sin oferta}

La tesis dice: \emph{``Las celdas $(j,a,w)$ en las que casi no hay oferta tienen precios sombra con
poco significado económico, así que reportamos solo estadísticas ponderadas por la masa de coches.''}
Aquí se desarrolla.

\subsection{Qué celdas son}

Hay celdas $(j,a,w)$ a las que casi ningún coche llega. Por ejemplo, $w$ muy lejos de 0 en edades
jóvenes: un coche nuevo nace con $w=0$ y en uno o dos años el desgaste no lo puede llevar al extremo
del grid. Si la transición de $w$ fuera la normal discretizada sin más, esas celdas tendrían masa
exactamente 0 (o $10^{-300}$). El código pone un piso: $F\leftarrow(1-\varepsilon_F)F+\varepsilon_F/W$ con
$\varepsilon_F=10^{-10}$, para que ninguna celda tenga masa cero y $\log S$ esté definido. La masa
de esas celdas queda del orden de $\varepsilon_F/W\approx10^{-12}$: viene del piso numérico, no de la economía.

\subsection{Por qué su precio sale alto}

En el logit, la demanda de una celda nunca es cero: cualquier comprador elige el coche $(j,a,w)$
con probabilidad positiva, $\Pi(y\mid i,b)\propto\exp\big(B(y)/\sigma_w\big)$, con $B(y)$ que
incluye $-\mu_\tau P_{jay}$. Al precio ``natural'' (el de sus vecinos) la demanda de la celda es del
orden de la de sus vecinos, por ejemplo $10^{-4}$, mientras que la oferta es $10^{-12}$. Para vaciar
el mercado, el precio tiene que subir hasta que la probabilidad de comprarla caiga por el mismo
factor. Como la oferta está acotada por la masa ($S\le q$), casi todo el ajuste lo hace la demanda,
y en el nido de $y$ esto requiere aproximadamente
$$\frac{\mu_\tau\,\Delta P}{\sigma_w}\approx\log\frac{D_{\text{natural}}}{S}
\quad\Longrightarrow\quad
\Delta P\approx\frac{\sigma_w}{\mu_\tau}\log\frac{D_{\text{natural}}}{S}.$$
Con $\sigma_w=0.1$, $\mu\approx0.1$ y una razón de $10^{8}$, $\Delta P\approx18$ mil DKK por encima de lo que
valdría un coche parecido. (Es un orden de magnitud ilustrativo: el ajuste real reparte el cambio entre
los niveles del árbol y entre tipos.)

\subsection{Por qué ``poco significado económico''}

\begin{enumerate}
  \item \textbf{Lo fija un artefacto numérico.} El precio depende de $\log(1/S)$, y $S$ depende de
  $\varepsilon_F$. Si $\varepsilon_F$ bajara de $10^{-10}$ a $10^{-20}$, con los mismos
  parámetros el precio subiría unos $\frac{\sigma_w}{\mu}\log(10^{10})\approx23$ mil DKK más. Cuando $\varepsilon_F\to0$,
  el precio de equilibrio de esa celda diverge: con demanda logit no existe un precio finito que haga
  la demanda igual a cero. Un precio que cambia al mover una constante numérica no dice nada sobre
  el mercado.
  \item \textbf{Es un precio sombra, no un precio de transacción.} En el sentido de Hicks (``precio
  virtual''), es el precio al que la demanda de un bien que no existe sería igual a su oferta nula.
  Con una demanda con precio de corte (choke price) sería el precio de reserva del comprador marginal.
  Con logit ese precio es infinito y el piso solo lo vuelve finito. Prácticamente nadie compra ni vende
  a ese precio: se intercambian $\sim10^{-12}$ coches.
  \item \textbf{Depende de la cola de los choques.} Que un coche que nadie tiene ``se demande'' un
  poco es consecuencia de los choques de valor extremo, que dan probabilidad positiva a cualquier
  alternativa. Es un supuesto de conveniencia, no una predicción sobre esas celdas.
  \item \textbf{No afecta al resto del equilibrio.} Esos precios entran al valor de los demás hogares
  solo con probabilidades minúsculas: la de llegar a esa celda ($\approx\varepsilon_F/W$) y la de
  comprarla (que es baja justo porque el precio es alto). Su efecto sobre $EV$, $q$ y los demás precios
  es del orden de $10^{-12}$ por decenas de utils, despreciable. Las celdas con masa, que son las que
  importan para la verosimilitud y los contrafactuales, tienen precios bien determinados.
  \item \textbf{Pueden verse raros.} Por lo anterior, un coche $(j,a,w)$ en una celda vacía puede salir
  más caro que el mismo coche con menos desgaste o menos edad. No es una contradicción del modelo: es
  el precio que raciona una demanda residual sobre una oferta casi nula.
\end{enumerate}

\subsection{Qué tan grande es el problema en el equilibrio base}

Con el equilibrio guardado en \texttt{output/modelo\_tesis/teoria/base} (régimen central, con el $R$
de esa corrida; residuo $\max|F|=3\times10^{-13}$), agregando los dos tipos:

\begin{center}
\begin{tabular}{lccc}
\toprule
Masa relativa de la celda & \% de celdas & Masa total & Transacciones\\
\midrule
$<10^{-12}$ & 6.1 & $2.5\times10^{-11}$ & $6.4\times10^{-11}$\\
$<10^{-8}$  & 10.2 & $2.3\times10^{-7}$ & $4.7\times10^{-7}$\\
$<10^{-6}$  & 21.6 & $1.2\times10^{-4}$ & $2.3\times10^{-4}$\\
$<10^{-4}$  & 59.0 & $3.8\times10^{-2}$ & $5.9\times10^{-2}$\\
\bottomrule
\end{tabular}
\end{center}

\noindent Masa y transacciones como fracción del total de coches activos y del total de ventas.
El síntoma más visible está en las edades jóvenes: un coche de un año de la marca ligera
($\overline P=142.3$) en una celda de $w$ extremo y vacía tiene precio de hasta 299.7 mil DKK, y uno de la
marca pesada ($\overline P=275.8$), hasta 485.8. Es decir, el usado vacío ``cuesta'' el doble que el nuevo.
En las celdas con masa relativa mayor a $10^{-6}$, el rango de un año es $[132.9,168.0]$ y $[245.3,286.9]$,
y el promedio ponderado por masa es 140.7 y 257.3. En un panel de $90{,}000$ transiciones, las celdas
con masa relativa menor a $10^{-8}$ aportan en promedio $90{,}000\times4.7\times10^{-7}\approx0.04$
observaciones por réplica.

Un fenómeno distinto, que no viene del piso: en las edades más viejas ($a=20$ y $a=24$), incluso en
celdas con masa, el precio no es monótono en $w$ (por ejemplo, 17 de 59 diferencias consecutivas van
``al revés'' en $a=24$ de la marca ligera). Ahí el rango de precios es estrecho ($[19.4,25.6]$) porque el
coche está a uno o pocos años de la terminal y $w$ casi no cambia su valor. Conviene revisarlo por
separado.

\subsection{Consecuencia para lo que se reporta}

Un promedio simple sobre las $W$ celdas de $(j,a)$ daría el mismo peso a una celda con $10^{-2}$ de
masa que a una con $10^{-12}$, y los precios sombra contaminarían el promedio. Por eso
\texttt{by\_age} reporta
$$\bar P^{\text{stock}}_{ja}=\frac{\sum_w q_{jaw}P_{jaw}}{\sum_w q_{jaw}},\qquad
\bar P^{\text{trans}}_{ja}=\frac{\sum_w S_{jaw}P_{jaw}}{\sum_w S_{jaw}},$$
ponderados por el número de coches en circulación o por el de transacciones. En ambos, una celda
con masa $10^{-12}$ pesa $10^{-12}$, así que estos promedios solo reflejan los precios que de verdad
se pagan.

\paragraph{Redacción sugerida para la tesis.} ``Hay celdas $(j,a,w)$ a las que casi ningún coche
llega, como desgastes extremos en edades jóvenes, con oferta del orden del piso numérico de la transición
de $w$. Como la demanda logit nunca es cero, su precio de equilibrio es el que raciona una demanda
residual sobre una oferta casi nula: es alto, depende de ese piso y casi no se intercambia a él.
Por eso reportamos precios promedio ponderados por la masa de coches o de transacciones, en los que
esas celdas no pesan.''

\end{document}
